% Source: https://artofproblemsolving.com/community/c3h2029806_happy_pi_day_maybe_a_competition
% Stable identity: aops-topic-2029806-q06
The given ring has area
\[
\pi\cdot5^2-\pi\cdot2.5^2=25\pi-6.25\pi=\dfrac{75\pi}{4}\ \text{cm}^2.
\]

\begin{asy}
usepackage("amsmath");
size(3cm);
defaultpen(fontsize(10pt));
fill(circle((0,0),2),paleblue); fill(circle((0,0),1),white);
draw(circle((0,0),2),linewidth(0.9)); draw(circle((0,0),1),linewidth(0.9));
draw((0,0)--2*dir(30),red); label("$5$",1.6*dir(30),NW,red);
draw((0,0)--dir(200),blue); label("$2.5$",0.5*dir(200),N,blue);
label("\mbox{}",(min(currentpicture,user=true)+max(currentpicture,user=true))/2,black); // keeps the text after the figure black when LaTeX typesets the labels inline
\end{asy}

The new ring has twice this area: \(2\cdot\dfrac{75\pi}{4}=\boxed{\dfrac{75\pi}{2}}\) cm\(^2\). (Its radii are \(\sqrt2\) times as large, \(5\sqrt2\) and \(\dfrac{5\sqrt2}{2}\) cm.)
